\documentclass[11pt]{article}
\usepackage[margin=1in]{geometry}
\usepackage{amsmath,amssymb,amsthm}
\usepackage{booktabs,longtable,array}
\usepackage{microtype}
\usepackage[hidelinks]{hyperref}
\newtheorem{theorem}{Theorem}
\newtheorem{lemma}{Lemma}
\newtheorem{proposition}{Proposition}
\newcommand{\vTwo}{\nu_2}
\newcommand{\ZZ}{\mathbb{Z}}
\newcommand{\QQ}{\mathbb{Q}}
\newcommand{\PPair}[2]{P_{#1,#2}}

\title{A Reduced Finite Certificate for the Six-Step Optimum\\
within a Fixed Four-State Tagged-Lift Family}
\author{Erd\H{o}s Problem 193 research note}
\date{13 September 2026}

\begin{document}
\maketitle

\begin{abstract}
For the radix-four triangular Gaussian walk used in a six-step construction for
Erd\H{o}s Problem 193, we consider all four-state integer tagged lifts of the
fixed form
\[
  W_n=4Z_n+d_{j_n},\qquad H_n=16n+c_{j_n},
\]
with arbitrary $d_j\in\ZZ^2$ and positive adjacent height increments. We give
a transparent computer-assisted finite proof that any such lift satisfying the
required $2$-adic valuation identity uses at least six distinct adjacent step
vectors. The certificate is solver-independent: it uses only exact rational
and integer arithmetic plus a complete residue enumeration modulo $16$.
The exhaustive classification reduces the $1050=S(8,5)$ exact five-block
partitions to
\[
  1050=184+164+675+27,
\]
and the final $27$ affine families are represented by only eight state-rotation
orbits.
\end{abstract}

\section{Setup}
Let
\[
 a=(1,i,-i,1),\qquad q_0=1,\qquad q_{4n+r}=a_rq_n,
 \qquad Z_n=\sum_{t<n}q_t.
\]
Write $q_n=i^{j_n}$, where $j_n\in\{0,1,2,3\}$. The adjacent state transitions
that occur are
\[
E=\bigl((0,1),(0,2),(1,2),(1,3),(2,3),(2,0),(3,0),(3,1)\bigr).
\]
We number these edges $e_0,\dots,e_7$ in the displayed order. Put
\[
 u_0=(1,0),\quad u_1=(0,1),\quad u_2=(-1,0),\quad u_3=(0,-1).
\]
For integer tags $d_j\in\ZZ^2$ and $c_j\in\ZZ$, with the gauge $d_0=(0,0)$
and $c_0=0$, define
\[
 W_n=4Z_n+d_{j_n},\qquad H_n=16n+c_{j_n}.
\]
If $j_n=j$ and $j_{n+1}=k$, the adjacent three-dimensional step is
\begin{equation}\label{eq:edge-step}
 S_{j,k}=\bigl(4u_j+d_k-d_j,\;16+c_k-c_j\bigr)\in\ZZ^3.
\end{equation}
We call the lift \emph{admissible} if every adjacent height increment is
positive and
\begin{equation}\label{eq:valuation}
 \vTwo\!\left(|W_n-W_m|^2\right)=\vTwo(H_n-H_m)\qquad(m<n).
\end{equation}
For convenience, $\PPair{m}{n}$ denotes the equality in \eqref{eq:valuation}
for the pair $(m,n)$.

\begin{lemma}[Automatic height-tag bound]\label{lem:c-bound}
Every lift with positive adjacent height increments satisfies
\[
 |c_j|\le30\qquad(j=0,1,2,3).
\]
\end{lemma}
\begin{proof}
For every transition $j\to k$ in $E$, $16+c_k-c_j\ge1$, hence
$c_k-c_j\ge-15$. The directed transition graph has diameter $2$. Thus a path
of length at most $2$ from $0$ to $j$ gives $c_j\ge-30$, while a path from
$j$ to $0$ gives $c_j\le30$.
\end{proof}

\begin{lemma}[Reduction to exactly five blocks]\label{lem:five-block}
To exclude all admissible lifts using at most five distinct steps, it suffices
to exclude all partitions of $E$ into exactly five blocks.
\end{lemma}
\begin{proof}
A lift using $k\le5$ step values induces a partition of the eight edges into at
most five equality classes. If $k<5$, split equality classes until exactly five
blocks remain. Splitting removes equality constraints, so the original lift
remains a solution of the refined five-block system.
\end{proof}

There are $S(8,5)=1050$ such partitions. We encode a partition by its
restricted-growth string $(\ell_0,\dots,\ell_7)$, written compactly as eight
digits; equal digits mean equal steps. State rotation $j\mapsto j+1\pmod4$
acts by $e_r\mapsto e_{r+2\bmod8}$ followed by canonical relabeling.

\section{The reduced certificate}
After the gauge choice there are nine unknown tag coordinates,
\[
(d_{1x},d_{2x},d_{3x},d_{1y},d_{2y},d_{3y},c_1,c_2,c_3).
\]
For any five-block partition, the equal-step conditions from
\eqref{eq:edge-step} are affine linear equations over $\QQ$. Exact Gaussian
elimination gives the following exhaustive split.

\begin{proposition}[Exhaustive partition split]\label{prop:split}
The $1050$ exact five-block partitions divide as follows:
\[
\begin{array}{c|r|l}
\text{class} & \text{count} & \text{reason for exclusion}\\ \hline
\text{linearly inconsistent}&184&\text{equal-step system inconsistent over }\QQ\\
\text{rank }9\text{, nonintegral}&164&\text{unique rational tag solution is not integral}\\
\text{rank }9\text{, integral}&675&\PPair{0}{4}\text{ or }\PPair{3}{7}\text{ fails}\\
\text{rank }6&27&\text{complete mod-16 residue certificate}\\ \hline
\text{total}&1050&
\end{array}
\]
\end{proposition}

\subsection{Linear obstructions}
Among the $184$ inconsistent five-block systems, $114$ already coarsen one of
eight inconsistent exact six-block systems. The eight six-block systems form
two $C_4$-rotation orbits. The remaining $70$ new minimal five-block
obstructions form twenty rotation orbits. Thus the linear part requires only
$22$ representative types; they are listed in Appendix~\ref{app:linear}.

\subsection{Full-rank systems}
There are $839$ consistent rank-$9$ systems. Their rational solution is unique.
Exactly $164$ have a nonintegral coordinate and are impossible for integer tags.
The remaining $675$ are integer solutions. Direct exact evaluation of only two
necessary valuation identities gives
\[
\begin{array}{c|r}
\text{failure pattern}&\text{count}\\\hline
\PPair{0}{4}\text{ and }\PPair{3}{7}\text{ both fail}&397\\
\text{only }\PPair{0}{4}\text{ fails}&139\\
\text{only }\PPair{3}{7}\text{ fails}&139\\
\text{both survive}&0.
\end{array}
\]
Hence no full-rank partition can be admissible.

\subsection{The rank-six exceptional families}
The remaining $27$ systems have rank $6$ and hence three affine degrees of
freedom. Exact row reduction verifies in every case that their integer tag
families have the common form
\begin{equation}\label{eq:rank6-form}
 d^x=p_x+Xv,\qquad d^y=p_y+Yv,\qquad c=Cv,
\end{equation}
for one three-vector $v$ and integer parameters $X,Y,C$, with the appropriate
integrality congruences encoded by the affine system.

By Lemma~\ref{lem:c-bound}, all valid height tags are finite. For the pairs used
below, $m,n\in\{0,\dots,7\}$, so
\[
0<H_n-H_m\le16\cdot7+60=172<256,
\]
and therefore $\vTwo(H_n-H_m)\le7$. Consequently the horizontal coordinates
need only be known modulo $16$. If $x,y$ are integers and
$\alpha=\vTwo(x)$, $\beta=\vTwo(y)$, then
\[
\vTwo(x^2+y^2)=
\begin{cases}
2\min(\alpha,\beta),&\alpha\ne\beta,\\
2\alpha+1,&\alpha=\beta,
\end{cases}
\]
so residues modulo $16$ determine the norm valuation whenever it can equal a
number at most $7$.

The checker enumerates all valid horizontal tag residues modulo $16$ and all
valid height tags, then imposes only
\begin{equation}\label{eq:four-pairs}
\PPair{2}{4},\qquad \PPair{1}{5},\qquad \PPair{1}{3},\qquad \PPair{0}{4}.
\end{equation}
Every one of the $27$ families is eliminated. Under state rotation they form
only eight orbits, listed in Appendix~\ref{app:rank6}.

\begin{theorem}[Six-step optimality in the fixed tagged-lift family]
Among all admissible integer tagged lifts
\[
W_n=4Z_n+d_{j_n},\qquad H_n=16n+c_{j_n},
\]
with positive adjacent height increments, the number of distinct adjacent
three-dimensional step vectors is at least six. The known six-step lift
therefore attains the minimum within this family.
\end{theorem}
\begin{proof}
By Lemma~\ref{lem:five-block}, a lift using at most five steps would satisfy one
of the $1050$ exact five-block equality systems. Proposition~\ref{prop:split}
and the rank-$6$ residue argument exclude every such system. Conversely, the
tags
\[
d_0=(0,0),\quad d_1=(-2,2),\quad d_2=(-2,-1),\quad d_3=(0,-3),
\]
\[
c=(0,8,5,13)
\]
produce the six step vectors
\[
\begin{split}
&(2,2,24),\quad (2,-1,21),\quad (0,1,13),\\
&(-2,-2,24),\quad (-2,1,11),\quad (0,-1,3),
\end{split}
\]
and satisfy the valuation identity for this construction. Hence the minimum is
exactly six.
\end{proof}

\section{Verification protocol}
The accompanying checker \texttt{verify\_erdos193\_reduced\_certificate.py}
uses only the Python standard library. It regenerates all set partitions,
forms the affine equal-step systems, performs exact row reduction with
\texttt{fractions.Fraction}, evaluates the rank-$9$ cases exactly, and performs
the complete rank-$6$ residue enumeration. No SMT solver, floating point
arithmetic, random sampling, or precomputed certificate is required. A separate
optional mode compares the regenerated $184$ linear contradictions against the
earlier certificate file. An independent audit run completed in under one
second on a desktop CPU.

The proof should therefore be regarded as a transparent finite
computer-assisted proof, not as a claim that the $1050$-case classification was
performed by hand.

\appendix
\section{Linear obstruction representatives}\label{app:linear}
The RGS strings below use the edge order
\[
(0,1),(0,2),(1,2),(1,3),(2,3),(2,0),(3,0),(3,1).
\]
The second column is the size of the representative's $C_4$ state-rotation
orbit within the indicated obstruction set.

\subsection*{Two exact six-block obstruction orbits}
\begin{center}
\begin{tabular}{cc}
\toprule
RGS representative & orbit size\\
\midrule
00122345 & 4 \\
01210345 & 4 \\
\bottomrule
\end{tabular}
\end{center}
Their eight members are already inconsistent. Exactly $114$ inconsistent
five-block systems are coarsenings of one of these eight systems.

\subsection*{Twenty new minimal exact five-block obstruction orbits}
\begin{center}
\begin{longtable}{cc}
\toprule
RGS representative & orbit size\\
\midrule
\endfirsthead
\toprule
RGS representative & orbit size\\
\midrule
\endhead
00012034 & 4 \\
00112341 & 4 \\
00112342 & 4 \\
00121334 & 4 \\
00121341 & 4 \\
00121342 & 4 \\
00123314 & 2 \\
00123441 & 4 \\
00123442 & 4 \\
00123443 & 4 \\
01022034 & 4 \\
01022134 & 4 \\
01023432 & 2 \\
01102340 & 4 \\
01120334 & 2 \\
01122034 & 4 \\
01201343 & 4 \\
01203142 & 4 \\
01230124 & 2 \\
01230143 & 2 \\
\bottomrule
\end{longtable}
\end{center}
The orbit-member counts sum to $70$.

\section{Rank-six orbit representatives}\label{app:rank6}
For each representative, the final column records the number of residue states
before any pair condition and after successively imposing
\[
\PPair{2}{4},\ \PPair{1}{5},\ \PPair{1}{3},\ \PPair{0}{4}.
\]
The initial state count equals $16^2$ times the number of valid height-tag
vectors.

\begin{center}
\scriptsize
\begin{tabular}{c c c c c}
\toprule
RGS & orbit & null vector $v$ & valid $c$ & residue-state progression\\
\midrule
00102343 & 4 & $(\frac{1}{2}, \frac{1}{2}, 1)$ & 23 & $5888\to2036\to0\to0\to0$ \\
00123240 & 4 & $(1, 1, 0)$ & 31 & $7936\to2724\to0\to0\to0$ \\
01012343 & 4 & $(\frac{1}{3}, \frac{2}{3}, 1)$ & 13 & $3328\to528\to2\to0\to0$ \\
01023241 & 4 & $(-1, -2, 1)$ & 13 & $3328\to528\to2\to0\to0$ \\
01123241 & 4 & $(0, -1, 1)$ & 23 & $5888\to2036\to2032\to2032\to0$ \\
01203041 & 4 & $(\frac{1}{2}, \frac{-1}{2}, 1)$ & 13 & $3328\to1156\to0\to0\to0$ \\
01213141 & 1 & $(1, 0, 1)$ & 31 & $7936\to0\to0\to0\to0$ \\
01213424 & 2 & $(1, 0, 1)$ & 31 & $7936\to0\to0\to0\to0$ \\
\bottomrule
\end{tabular}
\end{center}

\section{Initial base data}
The only base values used by the reduced local certificate are
\[
\begin{array}{c|c|c}
n&Z_n&j_n\\\hline
0&(0,0)&0\\
1&(1,0)&1\\
2&(1,1)&3\\
3&(1,0)&0\\
4&(2,0)&1\\
5&(2,1)&2\\
6&(1,1)&0\\
7&(2,1)&1
\end{array}
\]

\paragraph{Scope.}
The theorem above is an optimality statement only inside the fixed
$A=4$, $M=16$, four-state tagged-lift family. It does not assert that five
steps are impossible for Erd\H{o}s Problem 193 in general.

\end{document}
